\section{Design}\label{s:design}

This section provides a high-level overview of our design and system components. To execute an experiment, we first initialize Puffer's web and media servers in separate terminals. Then, we start a Mahimahi or CellReplay shell with a network trace we want to replay. Inside the emulated namespace, we run an automated script that launches a headless browser, logs in to Puffer's client, and starts a streaming session with the selected ABR algorithm. While Puffer logs application-layer metrics to its database, we also start a tcpdump process in another terminal to capture TCP packets. After the playback session ends, we extract the logged data as \emph{.CSV} and \emph{.PCAP} files. These data are later processed using Python scripts to generate visualizations and analyze the performance of ABR algorithms.

\textbf{Playback and Automation.}  
Playback is automated by a \emph{Puppeteer} script, a framework for \emph{Node.js}. Puppeteer provides a built-in headless browser, which we use to connect to Puffer's client. With this script, we reduce human errors by eliminating one-by-one command-line interactions. A timer of 10 minutes (which is the duration of the pre-recorded videos) is used for every test run to avoid looping. The timer sets the playback session to precisely 10 minutes and then closes the connection, so that each algorithm is tested for the same amount of time. For simplicity, we keep Puffer’s log-in page logic but include the log-in process inside the script by specifying the username and password, guaranteeing a setup that is less prone to human errors.

\textbf{ABR Configuration.}  
Each ABR algorithm requires configuration and setup through Puffer’s \emph{settings.yml} file. When this file is modified to change the ABR choice, we restart the media server before every run to reset the state and make sure that there is no carry-over between experiments. We keep the same configuration settings for every ABR to ensure fairness and accuracy in the results.

\textbf{Emulation Environment.}  
We initially chose \emph{Mahimahi} because of its popularity in academic ABR research, its deterministic trace-replay model, and its ability to easily integrate with Puffer. However, Mahimahi updates the bandwidth and delay only once per second and applies a fixed delay, which can be unrealistic in our time-sensitive setup. To address these limitations, we complement Mahimahi with \emph{CellReplay}, which introduces faster bandwidth and delay updates with more realistic 5G conditions. By comparing both emulators, we evaluate how the complexity of the emulator affects ABR performance.  
\begin{itemize}
    \item \textbf{Mahimahi} runs traces through \texttt{mm-link}, shaping up/down bandwidth and applying symmetric delay.  
    \item \textbf{CellReplay} runs traces through \texttt{mm-cell}, provides quicker updates for bandwidth and delay, supports asymmetric paths, and introduces base delays and latency offsets.
\end{itemize}

\textbf{Trace Collection.}  
We use real-world network traces in our experiments. Mahimahi includes AT\&T, T-Mobile, and Verizon LTE profiles in its repository. CellReplay provides high-resolution recordings from modern 5G networks, including three different T-Mobile and Verizon traces with varying signal strengths. To isolate emulator differences, we also run CellReplay’s 5G traces inside Mahimahi and check whether performance and rankings change for each ABR.

\textbf{Metric Collection.}  
Metrics such as SSIM, buffer occupancy, rebuffer time, RTT, and delivery rate are recorded by Puffer during playback in an InfluxDB database. Metrics such as retransmissions and TCP packets sent and received are monitored using \texttt{tcpdump}. We apply port filters to isolate streaming flows and use \texttt{tshark} to extract the percentage of retransmitted packets (ratio of retransmissions to total packets) and the timestamps of retransmissions (distribution of time between two successive retransmissions). These captures are processed by Python scripts for further analysis.

\textbf{Analysis and Visualization.}  
All plots, tables, and rankings in Section \ref{s:evaluation} are generated through a unified analysis pipeline that uses Python libraries such as \texttt{Matplotlib} and \texttt{Pandas}. We produce tables that summarize overall results and rank each ABR, scatter plots of mean SSIM vs. mean rebuffer time, RTT-over-time graphs showing how trace variability and emulator choice affect ABR behavior, and bar charts of average buffer time that reveal ranking changes across emulators.

%%% Local Variables:
%%% mode: latex
%%% TeX-master: "../thesis"
%%% End:
